\section{The random-sign IBLT}\label{sec:design}
We introduce the random-sign IBLT. The design multiplies each key by a random sign shared by both sides during insertion and deletion. The cell fields, field widths, and peeling process remain unchanged. We first give the encoding and decoding rules and prove ideal decoding equivalence with the standard IBLT and independence of the coding coefficients from the mapping. We then explain how the two-round protocol uses the lower bound $|b|$ to determine the second-round capacity.

\subsection{Encoding rules, decoding equivalence, and sign independence}
The two sides agree in advance on a sign hash $\varepsilon:\mathcal U\to\{-1,+1\}$ independent of the cell hash. Each sign is generated directly from the key and a shared seed, requiring no per-key storage. Algorithm~\ref{alg:encode} gives the complete encoding and decoding procedure. The lines marked $\star$ show all changes from the standard IBLT: applying the sign during insertion and deletion and recovering the direction on output. The number of cell fields, field widths, and peeling control flow remain unchanged. We describe encoding and decoding in turn.

\begin{algorithm}[!htbp]
\caption{Encoding and decoding of the random-sign IBLT. Lines marked $\star$ differ from the standard IBLT; field arithmetic is modulo $2^{64}$.}
\label{alg:encode}
\begin{algorithmic}[1]
\Procedure{Insert}{$x$}
\For{each cell $v$ of $x$}
\State $\mathrm{count}_v\gets\mathrm{count}_v+\varepsilon(x)$\Comment{$\star$ standard IBLT adds $1$}
\State $\mathrm{keySum}_v\gets\mathrm{keySum}_v+\varepsilon(x)\,x$\Comment{$\star$ standard IBLT adds $x$}
\State $\mathrm{checksum}_v\gets\mathrm{checksum}_v+\varepsilon(x)\,\mathrm{chk}(x)$\Comment{$\star$ standard IBLT adds $\mathrm{chk}(x)$}
\EndFor
\EndProcedure
\Procedure{Delete}{$x$}
\State as in \Call{Insert}{$x$}, subtract the same quantities from the three fields
\EndProcedure
\Procedure{Peel}{difference table}
\State add cells with count $\pm1$ to the candidate queue
\While{the candidate queue is nonempty}
\State remove cell $v$; if $C_v\ne\pm1$, process the next candidate
\State $x\gets C_v\,\mathrm{keySum}_v$
\If{$\mathrm{checksum}_v=C_v\,\mathrm{chk}(x)$ and the cells of $x$ include $v$}
\State output key $x$ and direction $C_v\,\varepsilon(x)$\Comment{$\star$ standard IBLT outputs $C_v$}
\State if the direction is $+1$, call \Call{Delete}{$x$}, otherwise call \Call{Insert}{$x$}
\State add cells whose count becomes $\pm1$ to the candidate queue
\EndIf
\EndWhile
\State \Return whether all three fields of every cell are zero
\EndProcedure
\end{algorithmic}
\end{algorithm}

Encoding follows lines 1--7 of Algorithm~\ref{alg:encode}. To insert a key $x$, we add $\varepsilon(x)$, $\varepsilon(x)\,x$, and $\varepsilon(x)\,\mathrm{chk}(x)$ to the three fields of each of its $k$ cells. Deletion subtracts the same quantities. We accumulate keySum and checksum modulo $2^{64}$. In the set model, these two fields also support the bitwise XOR implementation commonly used in the standard IBLT (Appendix~\ref{app:xor}). Our analysis and experiments use the additive form. Subtracting the two tables cell by cell cancels the identical contributions of each shared key in all three fields. A differing key $i$ has coding coefficient $s_i\varepsilon_i$ in the difference table, where $\varepsilon_i=\varepsilon(x_i)$. Thus, cell $v$ in the difference table has fields
\begin{equation}
C_v=\sum_{i:v\in S_i}s_i\varepsilon_i,\qquad
\mathrm{keySum}_v=\sum_{i:v\in S_i}s_i\varepsilon_i\,x_i,\qquad
\mathrm{checksum}_v=\sum_{i:v\in S_i}s_i\varepsilon_i\,\mathrm{chk}(x_i).\label{eq:fields}
\end{equation}
We treat $\varepsilon_i$ as iid uniform signs independent of the mapping $\mathcal G$.

Decoding follows lines 8--17 and uses the standard IBLT peeling process. A pure cell $v$ containing a single unrecovered differing key $x_i$ has count $C_v=s_i\varepsilon_i\in\{-1,+1\}$. Multiplying $\mathrm{keySum}_v$ by $C_v$ cancels the sign $s_i\varepsilon_i$ and yields the candidate key $x_i$ (line 12). Once the candidate passes the checksum and cell-membership checks (line 13), we multiply $C_v$ by $\varepsilon(x_i)$. Since $\varepsilon_i^2=1$, this recovers the original set direction $s_i=C_v\,\varepsilon(x_i)$ (line 14). We then call \textsc{Delete} or \textsc{Insert} according to the recovered direction $s_i$ (line 15). Both operations apply $\varepsilon(x_i)$ again in all $k$ cells of $x_i$. The contribution subtracted from each field is therefore exactly the key's contribution multiplied by its coding coefficient $s_i\varepsilon_i=C_v$. After each peeling step, the three fields still equal the encoded sum of the unrecovered keys. As in the standard IBLT, peeling succeeds when all three fields of every cell are zero (line 17).

Figure~\ref{fig:structure} compares the standard and random-sign IBLTs under the same mapping in a small one-sided example ($s_1=\dots=s_4=+1$). Both encodings have the same pure cells: the degree-$1$ cells 1, 4, 7, and 8, shaded in the figure. All have counts in $\{\pm1\}$. In pure cell 4, for example, the random sign gives $C_4=-1$. Multiplying by $\varepsilon_2=-1$ recovers the true direction of $x_2$ as $s_2=C_4\varepsilon_2=+1$. The two encodings follow identical peeling processes. Their sums of squared counts differ substantially. The single-key diagonal terms sum to $kd=12$ in both encodings. In the shared cells (cells 2, 3, 5, and 6), the standard IBLT accumulates contributions with the same sign. Each shared cell has count $2$, all cross terms are positive, and $\sum_vC_v^2=20$. In the random-sign IBLT, contributions along solid lines ($\varepsilon_i=+1$) cancel those along dashed lines ($\varepsilon_i=-1$) in the shared cells. Each shared cell in this example contains exactly one positive and one negative key. All cross terms are negative, the counts cancel to $0$, and $\sum_vC_v^2=4$. The next two propositions generalize the decoding equivalence and sign independence illustrated in Figure~\ref{fig:structure}.

\begin{figure}[!htbp]
\centering\includegraphics[width=.92\linewidth]{../figure/fig-structure.pdf}
\caption{The standard IBLT and the random-sign IBLT under the same mapping ($k=3$, $M=8$, one-sided difference $d=4$). The shaded cells have count $\pm1$; in (b), solid and dashed lines correspond to $\varepsilon_i=+1$ and $\varepsilon_i=-1$, respectively.}
\label{fig:structure}
\end{figure}

Random signs allow a degree-three cell with coding coefficients $(+1,+1,-1)$ to have count $+1$. This cannot occur in the standard IBLT under a one-sided difference. The candidate key read from such a cell fails the checksum and cell-membership checks, so peeling skips the cell until it contains a single key. Sign cancellation thus preserves correct key deletion and the stopping condition. Proposition~\ref{prop:decode} formalizes this property.

\begin{proposition}[Ideal decoding equivalence]\label{prop:decode}
Assume the pure-cell test has no false positives and field arithmetic does not overflow. At every intermediate state, the random-sign IBLT can remove exactly the keys removable by unsigned degree-one peeling. Under the same cell hash, the random-sign and standard IBLTs therefore leave the same 2-core. The failure event $F$ depends only on the unsigned mapping $\mathcal G$. The decoding threshold, failure probability, and recovered result are identical to those of the standard IBLT.
\end{proposition}

\begin{proof}
We induct on the number of peeling steps. By the field invariant, every degree-one cell has count $\pm1$ and passes the check. Ideal purity excludes cells of any other degree, so the two processes can follow the same cell-selection order. The full proof is in Appendix~\ref{app:proof-decode}.
\end{proof}

In an implementation, the finite 64-bit checksum gives a nonzero false-acceptance probability. Under an ideal hash, each check of a non-pure cell falsely accepts with probability at most $2^{-64}$. Section~\ref{sec:protocol} accounts for this term separately in the risk upper bound for the two-round protocol. Section~\ref{sec:exp-decode} verifies the implemented decoding equivalence in five hundred million paired trials with real IBLTs. Table~\ref{tab:decode-paired} in the appendix summarizes success-failure agreement and recovery checks for each configuration. Section~\ref{sec:exp-cost} reports decoding time.

Random signs retain their distribution after conditioning on failure. Failure selects mappings with large overlap, so conditioning usually changes the distribution of random variables associated with the mapping. This selection causes the bias of the standard IBLT under a one-sided difference. By Proposition~\ref{prop:decode}, failure depends only on the mapping. The signs come from a sign hash independent of the cell hash and remain iid uniform after we fix the mapping. Proposition~\ref{prop:composition} states this property precisely.

\begin{proposition}[Signs are independent of the mapping]\label{prop:composition}
For any fixed direction vector $(s_i)$, the coding coefficients $(s_i\varepsilon_i)$ are iid uniform signs independent of the mapping $\mathcal G$. Thus, failure $F$ is independent of the coding coefficients. Given $F$, or even the entire $\mathcal G$, the coding coefficients remain iid uniform signs. Statistics constructed from $(C_v)$, $\mathcal G$, and $F$ have the same joint distribution for all sign compositions $(d_+,d_-)$.
\end{proposition}

\begin{proof}
For fixed $(s_i)$, the map $(\varepsilon_i)\mapsto(s_i\varepsilon_i)$ is a bijection on $\{\pm1\}^d$ that preserves the uniform distribution. Since $(\varepsilon_i)$ is independent of $\mathcal G$, $(s_i\varepsilon_i)$ is an iid uniform sign sequence independent of $\mathcal G$. The joint distribution of $(s_i\varepsilon_i,\mathcal G)$ therefore does not depend on $(s_i)$. By Proposition~\ref{prop:decode}, $F$ is a function of $\mathcal G$, so conditioning on $F$ or $\mathcal G$ preserves the distribution of the coding coefficients. By \eqref{eq:fields}, $(C_v)$ is determined by $(s_i\varepsilon_i)$ and $\mathcal G$. Statistics constructed from $(C_v)$, $\mathcal G$, and $F$ thus have a joint distribution independent of $(s_i)$ and identical for all $(d_+,d_-)$.
\end{proof}

Section~\ref{sec:exp-e2e} and Table~\ref{tab:vldb} in the appendix illustrate this property numerically in the protocol. For a fixed configuration and calibration factor, random-sign encoding gives similar conditional miss rates under one-sided and balanced differences. Figure~\ref{fig:encoding} compares the post-failure distributions of the standard and random-sign encodings.

Propositions~\ref{prop:decode} and \ref{prop:composition} underpin our analysis. Decoding success depends only on $\mathcal G$, and the estimator's randomness comes from the independent signs. We can therefore analyze the post-failure distribution by first fixing the mapping and then studying the quadratic form in iid signs on the overlap graph. The results hold for all sign compositions. The joint distribution of the coding coefficients $s_i\varepsilon_i$ is independent of the direction vector $(s_i)$, so we can treat $s_i\varepsilon_i$ and $\varepsilon_i$ identically in the analysis. We write both as $\varepsilon_i$ below.

\subsection{\texorpdfstring{The two-round protocol flow and the lower bound $|b|$}{The two-round protocol flow and the lower bound |b|}}\label{sec:flow}
In the two-round protocol, the receiver subtracts the tables cell by cell and runs peeling on the difference table. For the standard IBLT, the count sum $\frac1k\sum_vC_v=b=d_+-d_-$ directly gives the deterministic lower bound $d\ge|b|$. For the random-sign IBLT, $\frac1k\sum_vC_v=\sum_i\varepsilon_i$ depends only on the random signs and no longer encodes the set-size difference $b$. The sender therefore attaches the set size $|A|$ (8 bytes) to the first-round table, allowing the receiver to compute $b=|A|-|B|$ directly.

If first-round peeling succeeds, the protocol ends immediately. After a first-round failure, the receiver computes an upper bound $U$ from the raw counts $(C_v)$ of the difference table. It incorporates the known lower bound $|b|$ by taking $U_b=\max(U,|b|)$ (for a one-sided difference, $|b|=d$). The second-round table has size $m(U_b)=\lceil U_b/(0.8\alpha_c)\rceil$. We re-encode and decode using a fresh cell hash independent of the first round. Whenever $U_b\ge d$, the second-round load is at most $0.8\alpha_c$, leaving a margin below the peeling threshold. The second round therefore succeeds with high probability. The analysis applies to any constant target load below $\alpha_c$; we use 0.8. Sections \ref{sec:measurement} and \ref{sec:inversion} analyze the post-failure distribution of the sum of squared counts and give the rule for computing $U$.
